\documentclass[11pt]{article}

\usepackage[margin=1in]{geometry}
\usepackage[T1]{fontenc}
\usepackage[utf8]{inputenc}
\usepackage{lmodern}
\usepackage{microtype}
\usepackage[round,authoryear]{natbib}
\usepackage[hidelinks]{hyperref}
\usepackage{doi}

\setlength{\bibsep}{0.5\baselineskip}

\title{Torc: Durable and Resource-Aware Workflow Orchestration\\
  from Workstations to Slurm Clusters}

% Replace this placeholder with the final author list, affiliations, and ORCIDs.
\author{Author list and affiliations to be finalized}
\date{Draft: \today}

\begin{document}

\maketitle

\begin{abstract}
Torc is an open-source workflow orchestration system for command-oriented computational pipelines.
It supports workflows ranging from independent parameter sweeps to dependency-rich, multi-stage
campaigns with heterogeneous CPU, memory, GPU, node-count, and runtime requirements. Researchers can
describe workflows declaratively, run them locally, and execute the same workflow through workers
inside Slurm allocations. A REST service and SQLite database maintain durable workflow state, while
a unified command-line interface supports creation, execution, monitoring, diagnosis, selective
reruns, and recovery. Python and generated API clients provide programmatic access, and terminal and
web interfaces support interactive operation. Torc is intended for computational scientists who
need a practical path from workstation-scale experiments to persistent HPC campaigns.
\end{abstract}

\section{Statement of Need}

Computational studies commonly begin with shell scripts, local process pools, or scheduler job
arrays. These approaches can be sufficient for independent tasks, but scientific campaigns often
grow to include dependencies, heterogeneous resource requirements, long execution histories, and
failures that must be diagnosed and selectively rerun. Moving such a campaign to an HPC system can
force researchers to change how they describe work, manage state, submit jobs, inspect progress, and
recover from failures.

Torc addresses this transition with a common workflow and control model across local and Slurm
execution. Workflow definitions describe commands, dependencies, files, user data, resource
requirements, schedulers, and failure behavior. Torc persists this information and execution state
behind an HTTP API. Workers claim ready jobs according to available resources, execute them as local
processes or Slurm job steps, and return results and resource observations. This separation allows
the same workflow to be inspected and operated through a command-line interface, Python client,
terminal interface, web dashboard, or direct API calls.

The compact control plane is important to Torc's intended use. Its server uses SQLite and does not
require researchers to administer a separate database service for core workflow state. Torc is not
intended to replace every workflow language, distributed runtime, or HPC scheduler. It targets
command-oriented scientific campaigns that benefit from durable coordination, resource-aware
execution, and a consistent path between workstations and Slurm clusters.

\section{State of the Field}

Scientific workflow systems span several overlapping design spaces. Snakemake and Nextflow provide
mature dataflow-oriented workflow languages and broad execution portability
\citep{molder2021snakemake,ditommaso2017nextflow}. Pegasus provides planning, provenance, data
management, and recovery for large scientific workflows \citep{deelman2019pegasus}. Parsl provides
Python-native parallel scripting and HPC executors \citep{babuji2019parsl}. FireWorks and Balsam
provide persistent services and worker-based execution for high-throughput scientific campaigns
\citep{jain2015fireworks,salim2018balsam}. Pilot systems and hierarchical schedulers, including
RADICAL-Pilot and Flux, execute many tasks within larger resource allocations
\citep{merzky2018radicalpilot,ahn2020flux}.

Torc does not claim to originate DAG execution, data-derived dependencies, persistent workflow
state, or task execution inside allocations. Its software contribution is a compact integration of
these ideas for command-oriented scientific users. Declarative specifications, a unified CLI,
durable SQLite-backed state, resource-aware local and Slurm workers, operational monitoring,
recovery, and provenance share one API and execution model. This design favors straightforward
deployment and consistent campaign operation over the language ecosystems, portable standards, or
specialized distributed runtimes offered by other systems.

\section{Software Design}

Torc uses a client-server architecture. An asynchronous Rust server exposes a versioned REST API and
stores workflows, jobs, dependencies, files, user data, resource requirements, results, events, and
scheduler records in SQLite. Clients and workers communicate with the server over HTTP, allowing the
database to remain on server-local storage while workers execute on other nodes.

Dependencies may be declared directly between jobs or inferred from producer and consumer
relationships over files and JSON user data. During workflow initialization, Torc resolves these
relationships into a common dependency graph. Workers atomically claim ready jobs and transition
them to a pending state, preventing concurrent workers from allocating the same job. Resource-aware
claims consider available CPUs, memory, GPUs, node count, and remaining runtime. Slurm workers run
inside allocations and can launch multiple resource-constrained job steps, enabling fine-grained
workflow execution without submitting every task independently to the scheduler.

Torc persists results and resource observations for later inspection. It supports failure handlers,
selective workflow reinitialization, resource correction after memory or runtime failures, and
offline completion journals when a worker temporarily loses access to the server. Optional RO-Crate
generation \citep{soilandreyes2022rocrate} records workflow, job, input, output, and software
provenance. Workflows can also extend their graphs at runtime through transactional job spawning for
adaptive or iterative algorithms.

Users interact with these capabilities through declarative workflow files and a unified
\texttt{torc} command-line interface. The same API supports generated clients, a Python
orchestration layer, a terminal UI, a web dashboard, and an MCP server. This layered design allows
interactive and programmatic clients to share workflow semantics rather than independently
implementing workflow state transitions.

\section{Research Impact Statement}

% Replace this paragraph with two or more concrete, anonymized deployments. For each deployment,
% report the scientific use, workflow shape, approximate scale, resource heterogeneity, execution
% environment, and concrete benefit provided by Torc.
Torc has been used for computational simulation campaigns, multi-stage data-processing pipelines,
and workflows with heterogeneous CPU and GPU requirements. These deployments motivate its emphasis
on durable state, resource-aware Slurm execution, monitoring, and selective recovery. The final
manuscript will describe at least two anonymized deployments and identify their scientific purpose,
workflow structure, approximate scale, execution environment, and realized research or operational
impact. Descriptive impact claims will be distinguished from performance claims requiring released
measurements and methods.

\section{Availability and Reproducibility}

Torc is distributed under the BSD 3-Clause license. Source code, documentation, examples,
integration tests, and release artifacts are publicly available in the project repository. The
project uses automated tests and generated-client consistency checks across its Rust server and
clients. The version reviewed with this paper will be archived with a persistent DOI, and citation
metadata will be provided in the repository.

% Add the repository URL, archived release DOI, exact paper version, and installation command.

\section{AI Usage Disclosure}

Generative AI tools were used to assist with literature discovery, manuscript brainstorming, and
early drafting. The authors will verify the resulting prose, technical claims, references, and
comparisons against the software, primary literature, and released evidence before submission. The
authors retain responsibility for the final manuscript.

\section*{Acknowledgements}

% Add institutional acknowledgements, funding sources, contract language, and contributor credit.
Acknowledgements and funding information will be added before submission.

\bibliographystyle{plainnat}
\bibliography{paper}

\end{document}
